\iffalse
!!!!!!!!!!!!!!!!!!!!!!!!!!!!!!!!!!!!!
! Do not modify this file directly  !
! This file has been auto-generated !
!!!!!!!!!!!!!!!!!!!!!!!!!!!!!!!!!!!!!
\fi

%% ODER: format ==         = "\mathrel{==}"
%% ODER: format /=         = "\neq "
%
%
\makeatletter
\@ifundefined{lhs2tex.lhs2tex.sty.read}%
  {\@namedef{lhs2tex.lhs2tex.sty.read}{}%
   \newcommand\SkipToFmtEnd{}%
   \newcommand\EndFmtInput{}%
   \long\def\SkipToFmtEnd#1\EndFmtInput{}%
  }\SkipToFmtEnd

\newcommand\ReadOnlyOnce[1]{\@ifundefined{#1}{\@namedef{#1}{}}\SkipToFmtEnd}
\usepackage{newtxmath}
\usepackage{stmaryrd}
\DeclareFontFamily{OT1}{cmtex}{}
\DeclareFontShape{OT1}{cmtex}{m}{n}
  {<5><6><7><8>cmtex8
   <9>cmtex9
   <10><10.95><12><14.4><17.28><20.74><24.88>cmtex10}{}
\DeclareFontShape{OT1}{cmtex}{m}{it}
  {<-> ssub * cmtt/m/it}{}
\newcommand{\texfamily}{\fontfamily{cmtex}\selectfont}
\DeclareFontShape{OT1}{cmtt}{bx}{n}
  {<5><6><7><8>cmtt8
   <9>cmbtt9
   <10><10.95><12><14.4><17.28><20.74><24.88>cmbtt10}{}
\DeclareFontShape{OT1}{cmtex}{bx}{n}
  {<-> ssub * cmtt/bx/n}{}
\newcommand{\tex}[1]{\text{\texfamily#1}}	% NEU

\newcommand{\Sp}{\hskip.33334em\relax}


\newcommand{\Conid}[1]{\mathit{#1}}
\newcommand{\Varid}[1]{\mathit{#1}}
\newcommand{\anonymous}{\kern0.06em \vbox{\hrule\@width.5em}}
\newcommand{\plus}{\mathbin{+\!\!\!+}}
\newcommand{\bind}{\mathbin{>\!\!\!>\mkern-6.7mu=}}
\newcommand{\rbind}{\mathbin{=\mkern-6.7mu<\!\!\!<}}% suggested by Neil Mitchell
\newcommand{\sequ}{\mathbin{>\!\!\!>}}
\renewcommand{\leq}{\leqslant}
\renewcommand{\geq}{\geqslant}
\usepackage{polytable}

%mathindent has to be defined
\@ifundefined{mathindent}%
  {\newdimen\mathindent\mathindent\leftmargini}%
  {}%

\def\resethooks{%
  \global\let\SaveRestoreHook\empty
  \global\let\ColumnHook\empty}
\newcommand*{\savecolumns}[1][default]%
  {\g@addto@macro\SaveRestoreHook{\savecolumns[#1]}}
\newcommand*{\restorecolumns}[1][default]%
  {\g@addto@macro\SaveRestoreHook{\restorecolumns[#1]}}
\newcommand*{\aligncolumn}[2]%
  {\g@addto@macro\ColumnHook{\column{#1}{#2}}}

\resethooks

\newcommand{\onelinecommentchars}{\quad-{}- }
\newcommand{\commentbeginchars}{\enskip\{-}
\newcommand{\commentendchars}{-\}\enskip}

\newcommand{\visiblecomments}{%
  \let\onelinecomment=\onelinecommentchars
  \let\commentbegin=\commentbeginchars
  \let\commentend=\commentendchars}

\newcommand{\invisiblecomments}{%
  \let\onelinecomment=\empty
  \let\commentbegin=\empty
  \let\commentend=\empty}

\visiblecomments

\newlength{\blanklineskip}
\setlength{\blanklineskip}{0.66084ex}

\newcommand{\hsindent}[1]{\quad}% default is fixed indentation
\let\hspre\empty
\let\hspost\empty
\newcommand{\NB}{\textbf{NB}}
\newcommand{\Todo}[1]{$\langle$\textbf{To do:}~#1$\rangle$}

\EndFmtInput
\makeatother
%
%
%
%
%
%
% This package provides two environments suitable to take the place
% of hscode, called "plainhscode" and "arrayhscode". 
%
% The plain environment surrounds each code block by vertical space,
% and it uses \abovedisplayskip and \belowdisplayskip to get spacing
% similar to formulas. Note that if these dimensions are changed,
% the spacing around displayed math formulas changes as well.
% All code is indented using \leftskip.
%
% Changed 19.08.2004 to reflect changes in colorcode. Should work with
% CodeGroup.sty.
%
\ReadOnlyOnce{polycode.fmt}%
\makeatletter

\newcommand{\hsnewpar}[1]%
  {{\parskip=0pt\parindent=0pt\par\vskip #1\noindent}}

% can be used, for instance, to redefine the code size, by setting the
% command to \small or something alike
\newcommand{\hscodestyle}{}

% The command \sethscode can be used to switch the code formatting
% behaviour by mapping the hscode environment in the subst directive
% to a new LaTeX environment.

\newcommand{\sethscode}[1]%
  {\expandafter\let\expandafter\hscode\csname #1\endcsname
   \expandafter\let\expandafter\endhscode\csname end#1\endcsname}

% "compatibility" mode restores the non-polycode.fmt layout.

\newenvironment{compathscode}%
  {\par\noindent
   \advance\leftskip\mathindent
   \hscodestyle
   \let\\=\@normalcr
   \let\hspre\(\let\hspost\)%
   \pboxed}%
  {\endpboxed\)%
   \par\noindent
   \ignorespacesafterend}

\newcommand{\compaths}{\sethscode{compathscode}}

% "plain" mode is the proposed default.
% It should now work with \centering.
% This required some changes. The old version
% is still available for reference as oldplainhscode.

\newenvironment{plainhscode}%
  {\hsnewpar\abovedisplayskip
   \advance\leftskip\mathindent
   \hscodestyle
   \let\hspre\(\let\hspost\)%
   \pboxed}%
  {\endpboxed%
   \hsnewpar\belowdisplayskip
   \ignorespacesafterend}

\newenvironment{oldplainhscode}%
  {\hsnewpar\abovedisplayskip
   \advance\leftskip\mathindent
   \hscodestyle
   \let\\=\@normalcr
   \(\pboxed}%
  {\endpboxed\)%
   \hsnewpar\belowdisplayskip
   \ignorespacesafterend}

% Here, we make plainhscode the default environment.

\newcommand{\plainhs}{\sethscode{plainhscode}}
\newcommand{\oldplainhs}{\sethscode{oldplainhscode}}
\plainhs

% The arrayhscode is like plain, but makes use of polytable's
% parray environment which disallows page breaks in code blocks.

\newenvironment{arrayhscode}%
  {\hsnewpar\abovedisplayskip
   \advance\leftskip\mathindent
   \hscodestyle
   \let\\=\@normalcr
   \(\parray}%
  {\endparray\)%
   \hsnewpar\belowdisplayskip
   \ignorespacesafterend}

\newcommand{\arrayhs}{\sethscode{arrayhscode}}

% The mathhscode environment also makes use of polytable's parray 
% environment. It is supposed to be used only inside math mode 
% (I used it to typeset the type rules in my thesis).

\newenvironment{mathhscode}%
  {\parray}{\endparray}

\newcommand{\mathhs}{\sethscode{mathhscode}}

% texths is similar to mathhs, but works in text mode.

\newenvironment{texthscode}%
  {\(\parray}{\endparray\)}

\newcommand{\texths}{\sethscode{texthscode}}

% The framed environment places code in a framed box.

\def\codeframewidth{\arrayrulewidth}
\RequirePackage{calc}

\newenvironment{framedhscode}%
  {\parskip=\abovedisplayskip\par\noindent
   \hscodestyle
   \arrayrulewidth=\codeframewidth
   \tabular{@{}|p{\linewidth-2\arraycolsep-2\arrayrulewidth-2pt}|@{}}%
   \hline\framedhslinecorrect\\{-1.5ex}%
   \let\endoflinesave=\\
   \let\\=\@normalcr
   \(\pboxed}%
  {\endpboxed\)%
   \framedhslinecorrect\endoflinesave{.5ex}\hline
   \endtabular
   \parskip=\belowdisplayskip\par\noindent
   \ignorespacesafterend}

\newcommand{\framedhslinecorrect}[2]%
  {#1[#2]}

\newcommand{\framedhs}{\sethscode{framedhscode}}

% The inlinehscode environment is an experimental environment
% that can be used to typeset displayed code inline.

\newenvironment{inlinehscode}%
  {\(\def\column##1##2{}%
   \let\>\undefined\let\<\undefined\let\\\undefined
   \newcommand\>[1][]{}\newcommand\<[1][]{}\newcommand\\[1][]{}%
   \def\fromto##1##2##3{##3}%
   \def\nextline{}}{\) }%

\newcommand{\inlinehs}{\sethscode{inlinehscode}}

% The joincode environment is a separate environment that
% can be used to surround and thereby connect multiple code
% blocks.

\newenvironment{joincode}%
  {\let\orighscode=\hscode
   \let\origendhscode=\endhscode
   \def\endhscode{\def\hscode{\endgroup\def\@currenvir{hscode}\\}\begingroup}
   %\let\SaveRestoreHook=\empty
   %\let\ColumnHook=\empty
   %\let\resethooks=\empty
   \orighscode\def\hscode{\endgroup\def\@currenvir{hscode}}}%
  {\origendhscode
   \global\let\hscode=\orighscode
   \global\let\endhscode=\origendhscode}%

\makeatother
\EndFmtInput
%
%


%
%
%
%
%
\ReadOnlyOnce{lineno.fmt}%

\usepackage{etoolbox}

\makeatletter

\newtoggle{lineno@column}
\newtoggle{lineno@print}
\newcounter{lineno@num}
\newlength\linenowidth

\newcommand\linenoFormat{\footnotesize}

\newcommand\linenoSetup{%
  \iftoggle{lineno@column}{%
    \column{lineno}{@{}r}%
    \column{B}{}
  }{}%
}

\newcommand\linenoStart{%
  \refstepcounter{lineno@num}%
  \iftoggle{lineno@print}{%
    \fromto{lineno}{B}{\makebox[\linenowidth][r]{\linenoFormat\arabic{lineno@num}}}%
  }{}%
}

\newcommand\linenoon{%
  \toggletrue{lineno@column}%
  \toggletrue{lineno@print}%
}

\newcommand\linenooff{%
  \togglefalse{lineno@column}%
  \togglefalse{lineno@print}%
}

\newcommand\linenoskip{%
  \toggletrue{lineno@column}%
  \togglefalse{lineno@print}%
}

\newcommand\linenoreset[1][1]{%
  \defcounter{lineno@num}{#1 - 1}%
}

\EndFmtInput












\section{Motivation}
\label{sec:motivation}

Context-free session types (CFSTs) extend regular session types with a general
sequential composition operator. A type $S_1;S_2$ describes a session that first
behaves as $S_1$ and then continues as $S_2$.  Composition is associative with
unit \ensuremath{\CodeTy{Skip}}, the inactive session.  Thus \ensuremath{\CodeTy{Skip}} and \ensuremath{\mathop{;}} form a monoid up to type
equality. The types \ensuremath{\CodeTy{Close}} and \ensuremath{\CodeTy{Wait}} indicate (active and passive) termination
of a session.

Sequential composition also interacts with choice in the expected way.
Appending a continuation to a choice appends it to every branch:
$(\oplus\{\ell_i:S_i\}_{i\in I});T =
\oplus\{\ell_i:S_i;T\}_{i\in I}$. The same law holds for external
choice. These equations make non-tail-recursive protocols
usable as session types: a continuation can be pushed through
the structure of a type until it reaches the leaves of a protocol.

As usual for session types, channels are linear resources. A channel
has to be used exactly once and every communication operation consumes
the channel at its current type and produces the channel at its
continuation type. Linearity is what makes the session type track the
state of the protocol, but it also forces programs to thread channels
explicitly from one operation to the next.

Additional remarks on design choices in this section, including why
some operations remain in resource-passing style, are given in the
supplement in \cref{sec:remarks}.

\subsection{Borrowing from Context-Free Session Types}
\label{sec:borr-from-cont}

\begin{figure}[tp]
  \TightCode
  \begin{subfigure}[t]{0.49\linewidth}
    \sethscode{tcolorboxcode}%
    \ApiOld
    \caption{CFST API}\label{lst:cfst-api}
  \end{subfigure}
  \begin{subfigure}[t]{0.49\linewidth}
    \ApiNew
    \caption{\thecalculus API}\label{lst:bcfst-api}
  \end{subfigure}
  \caption{Session type APIs}
  \label{fig:session-type-apis}
\end{figure}

Context-free session types enable writing elegant code for
serializing tree structures when communication is restricted to
primitive values.
% \footnote{We assume session-typed communication over a network
%   or between sandboxed processes, where only primitive data can be
%   passed, but no pointers.}
Consider an example from \citet{DBLP:conf/icfp/ThiemannV16} for serializing
binary trees of type \ensuremath{\Conid{Tree}}.
%
\DefnTree
%
\noindent Here is the corresponding session type \ensuremath{\Conid{TreeChan}}:
%
\DefnTreeChan
%
\noindent The type \ensuremath{\Conid{TreeChan}} is called context-free because it is
recursively defined and the recursion is not restricted to tail
positions. The \ensuremath{\CodeLbl{\Conid{Node}}} alternative is an example where the
\ensuremath{\Conid{TreeChan}} type refers two times in a row to itself. The
first occurrence of \ensuremath{\Conid{TreeChan}} in this
type is illegal in regular session type systems (e.g.,
\citet{DBLP:journals/jfp/GayV10,DBLP:conf/esop/LindleyM15,Gay_Vasconcelos_2025}) because
they syntactically restrict recursion to tail position. Context-free
session types lift this restriction as their composition operator $S_1;S_2$ connects
two sessions sequentially: the composed session comprises the
actions of $S_1$ followed by the actions of $S_2$.
% \footnote{One
%   might argue that CFSTs are close to Honda's original vision for
%   session types \cite{DBLP:conf/concur/Honda93}, which also features
%   symmetric composition.}

\Cref{lst:cfst-api} displays the API of CFSTs. The \ensuremath{\CodeOp{send}}, \ensuremath{\CodeOp{recv}}, and \ensuremath{\CodeOp{select}}
operations are standard for sending a message, receiving one, and selecting an
alternative in a branching session type. The \ensuremath{\CodeOp{branch}} operation maps a channel
offering different labeled alternatives to a variant type with the same labeled
alternatives. Standard pattern matching suffices to decode the chosen
alternative and continue. Given the types and the CFST API, we can implement the
serialization function:
%
\NumberedCode\CodeSendTree
%
\noindent However, there are several catches. First, the channel \ensuremath{\Varid{c}} is passed
explicitly throughout the code; a pattern we call \emph{resource-passing
  style}. This style forces the programmer to make the function polymorphic on
the continuation type of the channel, which is \ensuremath{\sigma } in this example. Second,
the \ensuremath{\CodeLbl{\Conid{Node}}} branch of \ensuremath{\Varid{sendTree}} relies on polymorphic recursion. To see
this, consider the type of \ensuremath{\Varid{c}_{2}} in
lines~\ref*{ln:sendTree/c2-def}--\ref*{ln:sendTree/c2-use}. To invoke \ensuremath{\Varid{sendTree}}
recursively, we need to instantiate the type argument \ensuremath{\sigma } with
\ensuremath{\Conid{TreeChan}\mathop{;}\sigma } rather than just \ensuremath{\sigma }, which is the telltale for polymorphic
recursion.

\BorrowCodeStyle

Our new language {\thecalculus} amends these problems. The type definitions
remain unchanged, but the implementation of \ensuremath{\Varid{sendTreeB}} relies on the session
API with borrowing shown in \cref{lst:bcfst-api}: the operations \ensuremath{\CodeOp{send}} and
\ensuremath{\CodeOp{recv}} adopt the borrowing style, but \ensuremath{\CodeOp{select}} and \ensuremath{\CodeOp{branch}} remain in
resource-passing style.
%% Reference is duplicated from start of the section.
% (see \cref{sec:remarks} for the reason).
To make the switch apparent, we present code using the \thecalculus{} API with
grey background and golden primitives.
%
\NumberedCode\CodeSendTreeB
%
\noindent What has changed? First, the type of \ensuremath{\Varid{sendTreeB}} is no longer
polymorphic over the continuation type. Instead, the channel argument \ensuremath{\Varid{c}} starts
out at type \ensuremath{\Conid{TreeChan}} without specifying a continuation. The return type is
\ensuremath{\CodeTy{Skip}}, which indicates an inactive channel.

Second, there is no explicit resource passing, except for \ensuremath{\CodeOp{select}}. Each use of
the borrow notation~\ensuremath{\mathord{\&\mkern-1mu}\Varid{c}} implicitly splits off a prefix of the session type
according to the function argument and updates \ensuremath{\Varid{c}} to the suffix. The comments
show the type of \ensuremath{\Varid{c}} after the preceding operation. Consider the type of \ensuremath{\Varid{c}}
after the \ensuremath{\CodeOp{select}} operation in line \ref*{ln:sendTreeB/select}, it starts
with~\ensuremath{\mathord{!}\Conid{Int}\mathop{;}\dots}. ``Taking the borrow''~\ensuremath{\mathord{\&\mkern-1mu}\Varid{c}} in the next line splits off the
borrowed type \ensuremath{\mathord{!}\Conid{Int}}---the argument type of \ensuremath{\CodeOp{send}\;\Varid{x}}---and updates \ensuremath{\Varid{c}} to the
type shown in line \ref*{ln:sendTreeB/send-x}. As the type of \ensuremath{\Varid{c}} is itself a
borrowed type, one could speak of a \emph{reborrow.}

Third, as there is no polymorphism, there is no polymorphic recursion, either.

Finally, the channel variable \ensuremath{\Varid{c}} can be used in a non-linear way, as long as
its uses are under borrows, as in \ensuremath{\mathord{\&\mkern-1mu}\Varid{c}}. The understanding is that the borrows
are taken and processed according to the evaluation order of the program. This
view does not abandon linearity. Rather, the type system internalizes the linear
use of the channel: each borrow consumes a distinct prefix of the current
session type and leaves the residual type for the next use of \ensuremath{\Varid{c}}.


\subsection{Local and Remote Borrowing}
\label{sec:local-remote-borr}

\Cref{sec:borr-from-cont} demonstrated \emph{local borrowing,} where the borrows
remain in the same process and are processed sequentially, following the
evaluation order.  In this local setting, operations on borrows are cheap:
borrowing merely names the next prefix of the same channel and leaves the suffix
for the surrounding computation. No other process can race with the
continuation, so the evaluation order of the program and the left-to-right order
enforced by typing already provide the required sequencing.

However, in some situations we want to send a borrow to
another process. As an example, suppose a server-side renderer streams
HTML nodes to a client instead of constructing a DOM tree in
memory first. Here is a protocol for such a renderer.
%
{\TightCode[\smallskipamount]\HtmlProtocolDefn}
%
\noindent The \ensuremath{\CodeLbl{\Conid{Elem}}} alternative first sends a tag name, and then a sequence
of children. The \ensuremath{\CodeLbl{\Conid{Child}}} alternative is context-free in the same sense as
the \ensuremath{\CodeLbl{\Conid{Node}}} alternative of \ensuremath{\Conid{TreeChan}}: rendering one child consumes an
\ensuremath{\Conid{HtmlChan}} prefix and then continues with the rest of the child stream.

\begin{Listing}
  \linenoon
  \linenoreset
  \TightCode
  \CodeRenderExample
  \caption{Rendering example}
  \label{lst:rendering-example}
\end{Listing}

Now consider the renderer presented in
\cref{lst:rendering-example}. \ensuremath{\Varid{renderUser}} emits a user panel whose first child
is computed by another process while the main process loads user messages, say,
from a database. We assume that \ensuremath{\Varid{renderProf}} and \ensuremath{\Varid{renderMsgArea}} render one HTML
node on a channel of type \ensuremath{\Conid{HtmlChan}\mathop{;}\CodeTy{Drop}} and \ensuremath{\Conid{HtmlChan}}, respectively. The type
\ensuremath{\CodeTy{Drop}} is a new session type that comes with a companion type \ensuremath{\CodeTy{Acq}}. They
indicate the placement of synchronization operations that ensure correct
handover of ownership from the borrowed session to the session's suffix.

The function \ensuremath{\Varid{renderUser}} starts an HTML element and then opens the child
stream. It passes a borrow of the first child position to a new process, which
renders the user's profile.  The continuation renders the remaining children,
here a messaging area, and then closes the child stream.

In this case, the borrow under the \ensuremath{\CodeOp{fork}} is a \emph{remote borrow}. Such a
borrow must be handled differently as operations on the borrow are no longer
necessarily ordered: the continuation must \emph{not} start emitting the second
child before the forked renderer has finished the first child. This
synchronization is expressed with the \ensuremath{\CodeOp{acquire}} operation, which works on a
channel of type \mbox{\ensuremath{\CodeTy{Acq}\mathop{;}\dots}.} It blocks the continuation until the borrowed
prefix has been completed by the forked process. Conversely, the forked process
calls \ensuremath{\CodeOp{drop}} on a channel of type \ensuremath{\CodeTy{Drop}} when it has finished working on the
prefix. Thus \ensuremath{\CodeOp{drop}} and \ensuremath{\CodeOp{acquire}} form the synchronization point between the two
processes; the underlying run-time mechanism is described
in~\cref{sec:impl-local-remote}.

This setup is also reflected in the types that result from a remote borrow. As
before, the comments show the type of \ensuremath{\Varid{c}} after the preceding
operation. Line~\ref*{ln:renderExample/fork} shows the type of the borrowed
segment passed to \ensuremath{\Varid{renderProfile}}. That session ends in \ensuremath{\CodeTy{Drop}}. After the \ensuremath{\CodeOp{fork}}
the type of \ensuremath{\Varid{c}} starts with \ensuremath{\CodeTy{Acq}}. The types \ensuremath{\CodeTy{Acq}} and \ensuremath{\CodeTy{Drop}} only arise from
remote borrows. The \emph{only} operation on a channel of type \ensuremath{\CodeTy{Acq}\mathop{;}\dots} is
\ensuremath{\CodeOp{acquire}} and, similarly, \ensuremath{\CodeOp{drop}} only applies to channels of type \ensuremath{\CodeTy{Drop}}.

Next, we discuss a direct operational semantics of local and remote borrowing
before we move on to a low level implementation of the two concepts in
\cref{sec:impl-local-remote}.


\subsection{A Direct Semantics for Channel Borrowing}
\label{sec:direct-semant-chann}

\begin{figure}
  \TightCode
  \setlength{\linenowidth}{1em}
  \begin{subfigure}[t]{.49\linewidth}
    \linenoon
    \linenoreset
    \CodeRenderRest
    \caption{Elaboration of \ensuremath{\Varid{renderRest}} to local splits.}
    \label{fig:rendering-example-elab/local}
  \end{subfigure}\hfill
  \begin{subfigure}[t]{.49\linewidth}
    \linenoon
    \linenoreset
    \CodeRenderExampleSnippetTrans
    \caption{Elaboration of \ensuremath{\Varid{renderUser}} to remote splits.}
    \label{fig:rendering-example-elab/remote}
  \end{subfigure}
  \caption{Snippets of the elaborated rendering example from
    \cref{lst:rendering-example}.}
  \label{fig:rendering-example-elab}
\end{figure}

\Cref{sec:borr-from-cont} already suggests an implementation strategy for local
borrowing. However, we start with an abstract specification because it is easier
to handle in a correctness proof. The specification works on an elaborated core
language, \thecalculus. Elaboration replaces borrowing with explicit
functional splitting operators. This translation is analogous to a translation
in previous work \citep{DBLP:journals/pacmpl/SaffrichS0V25}, which gives a formal
presentation along with a correctness proof. The novelty in our work is the
presence of \emph{two} operators, \ensuremath{\CodeOp{lsplit}} for local borrows and \ensuremath{\CodeOp{rsplit}} for
remote borrows. Borrows used in the same process elaborate to \ensuremath{\CodeOp{lsplit}} while
borrows used across a \ensuremath{\CodeOp{fork}} or sent to another process elaborate to
\ensuremath{\CodeOp{rsplit}}. \Cref{fig:rendering-example-elab} shows the (partial) elaboration of
\ensuremath{\Varid{renderRest}} and \ensuremath{\Varid{renderUser}} from the rendering example
(\cref{lst:rendering-example}), the former using local splits and the latter
using remote splits.

{\color{red}Our discussion starts with the local variant.} The translation of \ensuremath{\Varid{renderRest}} is
shown in \cref{fig:rendering-example-elab/local}. Execution proceeds by
repeatedly exposing the next local prefix of the channel
sequence. \Cref{fig:local-split-trace} shows a suffix of its reduction sequence.
The \mbox{\ensuremath{\boldsymbol{\nu}[\mkern1mu\relax \Varid{b}_{1}\mkern1mu][\mkern1mu\relax \Varid{b}_{2}\mkern1mu]\;\CodeKw{in}\;\dots}} operator is a channel restriction as known from
session calculi \citep{Gay_Vasconcelos_2025}. It introduces a communication
channel but with each compartment \ensuremath{[\mkern1mu\relax \Varid{b}_i\mkern1mu]} containing a \emph{bind group:} a
list of channel names and separator markers. The names bound by \ensuremath{\Varid{b}_{1}} refer to
one channel endpoint, the names in \ensuremath{\Varid{b}_{2}} to the other. The workings are best
illustrated with the example trace.

\begin{figure}
  \TightCode
  \begin{minipage}[t]{0.49\linewidth}
    \LocalSplitReduction
  \end{minipage}\hfill
  \begin{minipage}[t]{0.49\linewidth}
    \RemoteSplitReduction
  \end{minipage}\\
  \begin{subcaptiongroup}
    \parbox{.49\linewidth}{%
      \caption{Local split reduction suffix.}
      \label{fig:local-split-trace}}\hfill
    \parbox{.49\linewidth}{%
      \caption{Remote split and acquire reduction suffix.}
      \label{fig:remote-split-trace}}%
  \end{subcaptiongroup}
  \caption{Local and remote splitting in the suffixes of \ensuremath{\Varid{renderRest}} and
    \ensuremath{\Varid{renderUser}}. The processes \ensuremath{\Varid{P}_i} abbreviate client-side computation on the
    other endpoint.}
  \label{fig:local-split-reduction}
\end{figure}

To keep the trace short (\cref{fig:local-split-trace}), the reduction starts at
the local suffix in \ensuremath{\Varid{renderRest}}. The split on \ensuremath{\Varid{c}} following the \ensuremath{\CodeOp{select}}
\emph{replaces} the binder for \ensuremath{\Varid{c}} with two fresh channel names, \ensuremath{\Varid{c}_{1}} and \ensuremath{\Varid{c}_{2}}.
%
% All operations work on
% the channel names in a bind group from left to right.
%
% That is, both
% compartments of the \lstinline{nu} contain a flat list of channel names, and all
% operations work on these channel names from left to right. To ensure the correct
% sequencing, the semantics does not allow any action on \lstinline{c2} before
% \lstinline{c1} has been consumed.
%
This restriction is supported by typing: the \ensuremath{\CodeOp{lsplit}} operation returns an
\emph{ordered pair,} whose components must be processed and consumed in
left-to-right order.

% The |drop| operation (not used in this example) applied to some |c| just
% discards |c| if it is the first channel name in the compartment.

On the client side, in the processes \ensuremath{\Varid{P}_i} that hold the other channel end
\ensuremath{\Varid{d}}, splitting can be performed on \ensuremath{\Varid{d}} independently of the
splitting on \ensuremath{\Varid{c}}. The only constraint is that, in the end, when
\ensuremath{\CodeOp{close}} and \ensuremath{\CodeOp{wait}} calls match, the reduction requires that
only one channel name is left in each compartment. Under this condition, the
channel can be closed and eliminated. The figure does not show the transitions
of the {\color{red}client $P_d$,} which are independent.

For remote borrowing, we incorporate the \ensuremath{\CodeOp{acquire}} primitive and indicate a
remote split with the symbol ``\ensuremath{\mathop{\|}}'' in the compartment. As an example, we
consider the translation of lines
\ref*{ln:renderExample/transBegin}--\ref*{ln:renderExample/transEnd} from
\cref{lst:rendering-example}. The elaboration is shown in
\cref{fig:rendering-example-elab/remote}.

The translation inserts an \ensuremath{\CodeOp{rsplit}} because \ensuremath{\Varid{c}_{1}} and \ensuremath{\Varid{c}_{2}} are used in
different processes. The \ensuremath{\CodeOp{rsplit}} returns an ordinary (linear) pair of channels
because synchronization is made explicit using the \ensuremath{\CodeTy{Acq}} and \ensuremath{\CodeTy{Drop}}
types. Executing this code takes the steps shown in
\cref{fig:remote-split-trace} (ignoring the activity at the other end of the
channel). The bracketed step abbreviates the forked process running \ensuremath{\Varid{renderProfile}}
until its final \ensuremath{\CodeOp{drop}\;\Varid{c}_{1}}.  In summary, the \ensuremath{\CodeOp{rsplit}} replaces one channel name
by \ensuremath{[\mkern1mu\relax \Varid{c}_{1}\mathop{\|}\Varid{c}_{2}\mkern1mu]}, where ``\ensuremath{\mathop{\|}}'' marks a remote split. The \ensuremath{\CodeOp{drop}} is
only enabled if its argument name comes last before a ``\ensuremath{\mathop{\|}}''.  The
\ensuremath{\CodeOp{acquire}} is only enabled if the preceding list segment has been processed as in
\ensuremath{[\mkern1mu\relax \mathop{\|}\Varid{c}_{2}\mkern1mu]}. Executing the \ensuremath{\CodeOp{acquire}} in the last step removes the
exhausted segment and separator, yielding \ensuremath{[\mkern1mu\relax \Varid{c}_{2}\mkern1mu]}.


\subsection{Implementing Local and Remote Borrowing}
\label{sec:impl-local-remote}

We implement the high-level calculus by translation into a low-level run-time
session calculus with borrowing, \thertcalculus. The run-time calculus is a
concurrent lambda calculus with the usual session-style communication
primitives. However, \thertcalculus is untyped and it comes equipped with
asynchronous signaling primitives for implementing \ensuremath{\CodeOp{drop}}/\ensuremath{\CodeOp{acquire}}.  The direct
semantics from~\cref{sec:direct-semant-chann} records ordering constraints in
binding groups. The run-time translation realizes the same constraints with
ordinary channel names and optional synchronization flags that implement passing
of ownership.

\Cref{fig:borrow-lifecycle-msc} summarizes the difference between local
and remote borrowing in the implementation.

% Draws the figure with the local vs remote borrowing. Defines the following
% labels:
%
%   \label{fig:borrow-lifecycle-msc-local}
%   \label{fig:borrow-lifecycle-msc-remote}
%   \label{fig:borrow-lifecycle-msc}
% Make sure that any definitions we do in this file stay local.
\begingroup

\newlength\msgSelfWidth
\newlength\msgSelfHeight
\newlength\msgSelfOffset
\newlength\activationWidth
\newlength\firstStateOffset

\msgSelfWidth=0.9cm\relax
\msgSelfHeight=0.25cm\relax
\msgSelfOffset=0.15cm\relax
\activationWidth=0.14cm\relax
\firstStateOffset=0.08cm\relax

\newsavebox\SplitDrawingLocal
\newsavebox\SplitDrawingRemote

\newcommand\saveframedtikz[3][]{%
  \sbox{#2}{%
    \fbox{%
      \hspace{1mm}%
      \begin{tikzpicture}[#1]#3\end{tikzpicture}%
      \hspace{1mm}%
    }%
  }%
}

\tikzset{
  >=Stealth,
  head/.style={draw, minimum width=1.15cm,
    minimum height=0.42cm, inner sep=1pt, font=\scriptsize},
  state/.style={draw, fill=blue!5, align=center, inner sep=2pt,
    font=\scriptsize, text width=1.45cm, xshift = 2pt, anchor = west},
  flag/.style={draw=orange!65!black, fill=orange!14},
  inactive/.style={state, fill=black!12},
  lifeline/.style={dashed, black!55},
  activation/.style={draw, fill=black!8},
  msg/.style={->},
  create/.style={->, dashed},
  return/.style={->, dashed},
  lab/.style={font=\scriptsize},
  lab above/.style={lab,above,anchor=base,yshift=2.4pt},
  lab right/.style={lab,right,anchor=west,yshift=-.8pt},
}

\saveframedtikz\SplitDrawingLocal{
  \node[head] (p) at (0,0) {$P$};
  %\draw[lifeline] (p.south) -- ++(0,-5.3);

  % Define three points of the activation rectangle: upper right, lower
  % right and lower left.
  \path (p.south) ++(0, 0.2pt)
    ++( 0.5\activationWidth,  0) coordinate (act ur)
    ++( 0.0,                 -5) coordinate (act lr)
    ++(-1.0\activationWidth,  0) coordinate (act ll);

  % Draw the activation rectangle.
  \draw[activation] (act ur) rectangle (act ll);

  % Place nodes along the activation rectangle.
  \node[state] (c0) at ($(act ur)!.07!(act lr)$) {$c:S_1;S_2$};
  \node[state] (c1) at ($(act ur)!.35!(act lr)$) {$c_1:S_1$};
  \node[inactive, anchor = north west, yshift=.4pt] (c2 inactive)
    at ($(act ur)!(c1.south)!(act lr)$) {$c_2:S_2$\\inactive};
  \node[state] (c2) at ($(act ur)!.75!(act lr)$) {$c_2:S_2$};

  % Project (c1.north) onto the activation rectangle and draws the arrow
  % with the dimensions in \msgSelfWidth, \msgSelfHeight, \msgSelfOffset.
  \draw[msg] ($(act ur)!(c1.north)!(act lr)$)
    ++(0, \msgSelfOffset) ++(0, \msgSelfHeight)
    -- node[lab above] {$\CodeOp{lsplit}\,c$} ++(\msgSelfWidth,0)
    -- ++(0, -\msgSelfHeight)
    -- ++(-\msgSelfWidth, 0);

  % Repeat for the "use c1" arrow.
  \draw[msg] ($(act ur)!(c2.north)!(act lr)$)
    ++(0, \msgSelfOffset) ++(0, \msgSelfHeight)
    -- node[lab above] {use $c_1$} ++(\msgSelfWidth,0)
    -- ++(0, -\msgSelfHeight)
    -- ++(-\msgSelfWidth, 0);

  % Repeat for the "use c2" arrow. The arrow point is set 95% along the
  % activation rectangle.
  \draw[msg] ($(act ur)!.95!(act lr)$)
    ++(0, \msgSelfHeight)
    -- node[lab above] {use $c_2$} ++(\msgSelfWidth,0)
    -- ++(0, -\msgSelfHeight)
    -- ++(-\msgSelfWidth, 0);
}

\saveframedtikz\SplitDrawingRemote{
  \node[head] (p) {$P$};
  \node[head, flag, below right=0.37cm and 2.5cm of p] (z) {$\varphi\,z$};
  \node[head,       below right=0.61cm and 2.5cm of z] (q) {$Q$};

  % Coordinates for the upper edge of p's first activation.
  \path (p.south) ++(0, .2pt)
    + (-.5\activationWidth,  0) coordinate (p1 ul)
    ++( .5\activationWidth,  0) coordinate (p1 ur)
    ++(                  0, -1) coordinate (p er);

  % The initial c channel.
  \path (p1 ur) ++(0, -\firstStateOffset)
    node[state, anchor=north west] (c) {$c:S_1;S_2$};

  % p's lifeline
  \draw[lifeline] (p.south) -- ++(0, -5) coordinate (p end);

  % Define three points of p's first activation rectangle: upper right,
  % lower right and lower left.
  \path (p1 ur)
    ++( 0.0,                 -2) coordinate (p1 lr)
    ++(-1.0\activationWidth,  0) coordinate (p1 ll);
  \draw[activation] (p1 ur) rectangle (p1 ll);

  % q's activation rectangle
  \path ($(q.north)!(p end)!(q.south)$)
    ++(0, .5)
    +( .5\activationWidth, 0) coordinate (q lr)
    +(-.5\activationWidth, 0) coordinate (q ll);
  \path (q.south) +(.5\activationWidth, .2pt) coordinate (q ur);
  \draw[activation] (q ur) rectangle (q ll);

  % Define coordinates that define the right edge of z's activation
  % rectangle.
  \path (z.south) ++(0, .2pt)
    ++(0.5\activationWidth, 0) coordinate (z ur)
    ++(0, -1) coordinate (z er);

  % actions and states in Q
  \path (q ur) ++(0, -\firstStateOffset)
    node[state, anchor=north west] (c1 at q) {$c_1:S_1;\CodeTy{Drop}$};
  \draw[msg] ($(q ur)!(c1 at q.south)!(q lr)$)
    ++ (0, -\msgSelfOffset)
    -- ++(\msgSelfWidth,0)
    -- node[lab right] {use $c_1$} ++(0, -\msgSelfHeight)
    -- ++(-\msgSelfWidth, 0) coordinate (q use end);
  \path (q use end) ++(0, -\msgSelfOffset)
    node[state, anchor=north west] (c1 drop) {$c_1:\CodeTy{Drop}$};

  % the "drop" arrow
  \path ($(q ur)!(c1 drop.south)!(q lr)$)
    ++(-\activationWidth,-\msgSelfOffset)
    coordinate (drop origin);
  \draw[msg] (drop origin) -- node[lab above] {$\CodeOp{drop}\,c_1$}
    ($(z ur)!(drop origin)!(z er)$) coordinate (drop target);

  % States for the synchronization flag.
  \path (z ur) ++(0, -\firstStateOffset)
    node[state, flag, anchor = north west] {$z\mapsto\texttt{drop}$};
  \path (drop target) ++(0, -\msgSelfOffset)
    node[state, flag, anchor = north west] (z acq) {$z\mapsto\texttt{acq}$};

  % Define the remaining coordinates for z's activation rectangle.
  \path (z ur) ++(-\activationWidth, 0) coordinate (z ul);
  \path ($(z ur)!(z acq.south)!(z er)$)
    ++(0, -\msgSelfOffset) coordinate (z lr)
    ++(-\activationWidth, 0) coordinate (z ll);
  \draw[activation, flag] (z ur) rectangle (z ll);

  % Define three points of p's second activation rectangle: lower right,
  % lower left and upper right.
  \path (p end) +( .5\activationWidth, 0) coordinate (p2 lr)
                +(-.5\activationWidth, 0) coordinate (p2 ll);
  \coordinate (p2 ur) at ($(p2 lr)!(z ll)!(p1 lr)$);
  \draw[activation] (p2 ur) rectangle (p2 ll);

  % z's creation arrow
  \draw[create]
    ($(p1 ur)!(z.west)!(p1 lr)$) coordinate (rsplit origin)
    -- node[lab above, near end] {$\CodeOp{rsplit}\,c$}
       node[lab, below, yshift=2.2pt, near end] {create $z$} (z.west);

  % c1 and c2 immediately after the rsplit
  \path (rsplit origin) ++(0, -\msgSelfOffset)
    node[state, anchor = north west] (c1 at p) {$c_1:S_1;\CodeTy{Drop}$};
  \node[state, anchor = north west, yshift=.4pt] (c1 at p)
    at ($(p1 ur)!(c1 at p.south)!(p er)$) {$c_2:\mathrlap{\CodeTy{Acq};S_2}\hphantom{S_1;\CodeTy{Drop}}$};

  % Remember the coordinate where the fork arrow starts.
  \coordinate (fork q) at ($(p1 ur)!(q.west)!(p1 lr)$);

  % q's fork arrow
  \draw[create] (fork q) -- (q);
  \path (fork q) -- node[lab above, near end] {$\CodeOp{fork}$}
                    node[state, below=2pt, near end] {$c_1 : S_1;\CodeTy{Drop}$} (q);

  % the "acquire" arrow
  \draw[msg] (p1 lr) --
    node[lab above] {$\CodeOp{acquire}\,c_2$}
    ($(z ul)!(p1 lr)!(z ll)$);
  % the "return" arrow
  \draw[return] (z ll) -- node[lab above] {return} (p2 ur);

  \path (p2 ur) ++(0, -\msgSelfOffset)
    node[state, anchor=north west] (c2 ret) {$c_2:S_2$};
  \draw[msg] ($(p2 ur)!(c2 ret.south)!(p2 lr)$) ++(0,-\msgSelfOffset)
    -- ++(\msgSelfWidth,0)
    -- node[lab right] {use $c_2$} ++(0,-\msgSelfHeight)
    -- ++(-\msgSelfWidth,0);
}

\begin{figure}[t]
  \begin{subfigure}[t]{\wd\SplitDrawingLocal}
    \usebox\SplitDrawingLocal
    \caption{Local borrow}
    \label{fig:borrow-lifecycle-msc-local}
  \end{subfigure}\hfill
  \begin{subfigure}[t]{\wd\SplitDrawingRemote}
    \usebox\SplitDrawingRemote
    \caption{Remote borrow}
    \label{fig:borrow-lifecycle-msc-remote}
  \end{subfigure}
  \caption{Message-sequence view of local and remote borrows. Channel state is
    shown as annotations next to the relevant process or flag lifelines. A local
    split keeps the channel state with the original process lifeline; a remote
    split creates a flag lifeline $\varphi\,z$ and forks a process for the
    borrowed prefix.}
  \label{fig:borrow-lifecycle-msc}
\end{figure}

% Closes the group opened at the start of the file.
\endgroup

%%% Local Variables:
%%% mode: LaTeX
%%% TeX-master: "popl2027"
%%% End:


%
% Locally format "in" as a variable name
%

Specifically, a channel is represented by a triple \ensuremath{\Varid{in}\otimes \Varid{c}\otimes \Varid{out}}, where \ensuremath{\Varid{c}}
is a primitive channel name and \ensuremath{\Varid{in}} and \ensuremath{\Varid{out}} are optional names of
synchronization flags. The incoming flag \ensuremath{\Varid{in}} records whether the channel must
be acquired before it can be used. The outgoing flag \ensuremath{\Varid{out}} records whether the
channel must be dropped after use. A synchronization binder
{\color{red}$(\varphi z \mapsto \phi)P$} binds a flag name \ensuremath{\Varid{z}} in \ensuremath{\Conid{P}} and stores
one of two states in {\color{red}$\phi$}: \ensuremath{\CodeFlag{drop}} or \ensuremath{\CodeFlag{acq}}.  This flag \ensuremath{\Varid{z}}
implements the separator in \ensuremath{\boldsymbol{\nu}[\mkern1mu\relax \Varid{c}_{1}\mathop{\|}\Varid{c}_{2}\mkern1mu][\mkern1mu\relax \dots\mkern1mu]}. It is the outgoing flag of
the borrowed prefix and the incoming flag of the suffix.

An \ensuremath{\CodeOp{rsplit}} creates a flag in state \ensuremath{\CodeOp{drop}} and a freshly restricted channel sets
its flag to \ensuremath{\Varid{acq}}. Intuitively, \ensuremath{\CodeFlag{drop}} means that the left, borrowed part of a
remote split is still active, \ensuremath{\CodeFlag{acq}} means that it has been dropped and the
right, suffix part may acquire. This mechanism implements an ownership transfer
of the primitive channel from the borrowed part to the suffix.

Allocation creates a primitive channel pair and one synchronization flag for
each endpoint. The fresh endpoints carry incoming flags that are already in the
\ensuremath{\CodeFlag{acq}} state. Thus, an initial \ensuremath{\CodeOp{acquire}} is required to reduce each channel to
normal operation, removing the flag from the channel triple by replacing it with
\ensuremath{\mathord{*}}. Once these optional flag components are absent, communication is just
ordinary session communication: \ensuremath{\CodeOp{send}}/\ensuremath{\CodeOp{recv}}, \ensuremath{\CodeOp{select}}/\ensuremath{\CodeOp{branch}}, and
\ensuremath{\CodeOp{close}}/\ensuremath{\CodeOp{wait}} all operate directly on the primitive channel name \ensuremath{\Varid{c}}.

\Cref{fig:borrow-lifecycle-msc-local} shows the lifecycle of a local borrow. The
rule for \ensuremath{\CodeOp{lsplit}} returns two copies of the same run-time channel triple. No new
channel is allocated and no synchronization flag is introduced. This simple
approach has no overhead and is sound because the source type system has already
established the left-to-right use of the two pieces. The run-time representation
can therefore be imperative: the two pieces are aliases for the same underlying
channel object, and using the first piece advances that object to the state
where the second piece starts.

Remote splitting is the only place where new synchronization machinery is
introduced (\cref{fig:borrow-lifecycle-msc-remote}). Splitting a channel \ensuremath{\Varid{in}\otimes \Varid{c}\otimes \Varid{out}} remotely creates a fresh flag \ensuremath{\Varid{z}} in state \ensuremath{\CodeFlag{drop}} and returns two
channel triples: \ensuremath{\Varid{in}\otimes \Varid{c}\otimes \Varid{z}} for the borrowed prefix, and \ensuremath{\Varid{z}\otimes \Varid{c}\otimes \Varid{out}} for
the suffix. The prefix therefore carries an outgoing obligation to call \ensuremath{\CodeOp{drop}};
the suffix carries a matching incoming obligation to call \ensuremath{\CodeOp{acquire}}. Calling
\ensuremath{\CodeOp{drop}} changes the flag from \ensuremath{\CodeFlag{drop}} to \ensuremath{\CodeFlag{acq}} and returns immediately, so the
notification is asynchronous. A later \ensuremath{\CodeOp{acquire}} on the suffix can fire only
after this state change and removes the synchronization flag. Thus a remote
split pays for exactly the synchronization needed to coordinate the two
processes operating on the channel, whereas a local split pays nothing beyond
sharing the channel representation.

% End local formatting of "in" as an identifier.
%

%%% Local Variables:
%%% mode: latex
%%% TeX-master: "../popl2027.tex"
%%% End:
